% A.2 and A.4.  Authoring agent: GPT-6 Astra (OpenAI), run /root/a2_a4.
% This fragment uses only standard amsmath/amssymb/amsthm commands.

\section{A.2. The walking complex and the exact universal property}
\label{sec:A2}

We use homological degrees for complexes in the source and cohomological
degrees for complexes in the target.  Let $R$ be a commutative ring with
identity, including the zero ring.  Write
\[
 \mathcal P_+(R)=\operatorname{Ch}^{b}_{\geq0}
                 (\operatorname{Free}_{\mathrm{fg}}R).
\]
Its morphisms are chain maps; its distinguished short exact sequences are
those split in every degree.  The disk $D_i$, for $i\geq0$, has $R$ in
degrees $i+1,i$ and differential $1_R$.  The stalk $S_i$ has $R$ in degree
$i$.  The canonical map $D_i\to D_{i+1}$ is the identity from the upper
copy of the first disk to the lower copy of the second.  We denote this
map by $u_i$; then $u_{i+1}u_i=0$.  These are the manuscript's $C^i$ and
$\Psi^i$.  The nonnegative indexing matters: scan 49 says $n\geq0$, and
scan 52 uses the sequences $0\to\Psi^i\to C^i\to\Psi^{i+1}\to0$.

An $R$-linear additive functor $F:\mathcal P_+(R)\to\mathcal B$ is
\emph{relatively left exact} if it sends every distinguished sequence
$0\to A\to B\to E\to0$ to a kernel diagram
$F(A)\to F(B)\to F(E)$.  This definition uses ordinary categorical
kernels in $\mathcal B$.  We treat the stronger requirement that the
kernel inclusions be admissible separately below.  This distinction is
needed because the manuscript uses ``$\Sigma$-kernel'' before its exact
structure axioms have been completed.

\paragraph{A2.1 (Refuted, with corrected statement).}
\label{A2.1}
For every nonzero $R$, the category $\mathcal P_+(R)$ is not Frobenius.
It has enough projectives and enough injectives, but the stalk $S_0$ is
projective and is not injective.  Thus the identification in Lead L3 and
in the study's discussion of 137:49--56 is false for the category actually
defined there.  Its valid two-sided replacement is A2.6.

\begin{proof}
A chain map $S_0\to L$ is precisely a map $R\to L_0$.  Since the
distinguished sequences split in degree zero, $S_0$ is projective.
The distinguished sequence
\[
 0\longrightarrow S_0\xrightarrow{j}D_0
   \longrightarrow S_1\longrightarrow0
\]
does not split as a sequence of complexes.  Indeed, a chain map
$r:D_0\to S_0$ has degree-one component zero; the chain-map equation
forces its degree-zero component to be zero.  Therefore $rj$ cannot be
$1_{S_0}$ when $R\ne0$.  This proves noninjectivity.  For $R=\mathbf F_2$
the finite certificate enumerates both possible degree-zero scalars and
finds no retraction.

For completeness, the disks are projective and injective: there are
natural identifications
\[
 \operatorname{Hom}(D_i,L)=L_{i+1},\qquad
 \operatorname{Hom}(L,D_i)=\operatorname{Hom}_R(L_i,R).
\]
Both functors take degreewise split sequences to exact sequences.
Every $L$ admits a degreewise split epimorphism from
$S_0\otimes_R L_0\oplus\bigoplus_{i\geq1}D_{i-1}\otimes_RL_i$;
the map in degree $n$ sends $(x,y)$ to $x+\partial y$.
Its kernel in degree $n$ is isomorphic to $L_{n+1}$, so belongs to
$\mathcal P_+(R)$.  The injective embedding is constructed in A2.2.
All sums here are finite.  For the zero ring the category is zero and
the exceptional non-Frobenius assertion is, as stated, excluded.
\end{proof}

For a finite free module $V$ and an object $B$ of an $R$-linear additive
category, $V\otimes_R B$ denotes its finite copower.  Choosing a basis
identifies it with a finite direct sum of $B$; a change of basis acts by
the corresponding invertible matrix.  Consequently all the following
formulas are independent of bases, up to the unique comparison
isomorphisms induced by those matrices.  No tensor product between two
objects of $\mathcal B$ is assumed.

\paragraph{A2.2 (Proved; every commutative base ring).}
\label{A2.2}
Let $\mathcal B$ be an $R$-linear additive category.  Restriction
\[
 \rho:F\longmapsto \bigl(F(D_i),F(u_i)\bigr)_{i\geq0}
\]
from relatively left-exact functors to nonnegative cochain complexes in
$\mathcal B$ is fully faithful.  Its essential image consists exactly
of the complexes $C=(C^i,d^i)$ for which the following explicitly typed
kernel exists for every bounded nonnegative finite-free chain complex
$L=(L_i,\partial_i)$:
\begin{align*}
 \delta_L(C):\bigoplus_{i=0}^{m}L_i\otimes_R C^i
   &\longrightarrow \bigoplus_{i=0}^{m}L_i\otimes_R C^{i+1},\\
 (x_i)_{i=0}^{m}&\longmapsto
 \bigl((1\otimes d^i)x_i-(\partial_{i+1}\otimes1)x_{i+1}\bigr)_{i=0}^{m}.
\end{align*}
We set $L_{m+1}=0$ and $x_{m+1}=0$.  An inverse on the essential image is
\[
 T_C(L)=\ker\delta_L(C).
\]
Thus the answer for an arbitrary commutative ring, including any
principal ideal ring with zero divisors, is a family of finite-matrix
kernel conditions.  The matrices $\partial_i$ range over all finite
matrices satisfying $\partial_i\partial_{i+1}=0$.  When $\mathcal B$
has all kernels, every cochain complex is in the image.

\begin{proof}
For $b\in\mathcal B$, write $C_b^i=\operatorname{Hom}_{\mathcal B}(b,C^i)$.
Let $L^\vee$ be the cochain complex with
$(L^\vee)^i=\operatorname{Hom}_R(L_i,R)$ and differential
$\partial_{i+1}^{\vee}$.  By the finite-copower universal property,
\[
 \operatorname{Hom}_{\mathcal B}(b,T_C(L))
 =\operatorname{Hom}_{\operatorname{Ch}(R)}(L^\vee,C_b).
 \tag{A2.2a}
\]
This equality proves that the kernel defines a covariant additive
functor of $L$ and a covariant functor of $C$.  The right side is the
group of degree-preserving cochain maps, not a derived Hom group.
Dualizing a degreewise split short exact sequence in $\mathcal P_+(R)$
gives a short exact sequence of cochain complexes in the reverse order.
The contravariant Hom functor into $C_b$ is left exact.  Applying this
for every $b$ proves that $T_C$ is relatively left exact.  On $D_i$ a
cochain map $D_i^\vee\to C_b$ is uniquely specified by its component in
degree $i$.  Hence $T_C(D_i)=C^i$, naturally, with the stated
differential.  In particular $\rho T_C\cong C$.

We now prove that every relatively left-exact $F$ has the displayed
form, without imposing a Frobenius hypothesis.  Define a functor
\[
 I(L)=\bigoplus_{i\geq0}D_i\otimes_RL_i.
\]
In degree $n$ this is $L_n\oplus L_{n-1}$, where $L_{-1}=0$, and its
differential sends $(x,y)$ to $(y,0)$.  The chain map
\[
 j_L:L\longrightarrow I(L),\qquad x\longmapsto(x,\partial x)
\]
is a degreewise split monomorphism.  Its quotient $Q(L)$ has
$Q(L)_n=L_{n-1}$ and differential $-\partial$.  Explicitly the quotient
map is $q_L(x,y)=y-\partial x$.  Thus
\[
 0\longrightarrow L\xrightarrow{j_L}I(L)
      \xrightarrow{q_L}Q(L)\longrightarrow0
\]
is distinguished, as is the analogous sequence for $Q(L)$.
Relative left exactness implies
\[
 F(L)=\ker\bigl(F(I(L))\xrightarrow{F(j_{Q(L)}q_L)}F(I(Q(L)))\bigr).
\]
If $C=\rho F$, additivity identifies the displayed map with
$\delta_L(C)$: its disk-to-next-disk entries are $d^i$, and its
same-disk entries are $-\partial_{i+1}$.  This proves
$F\cong T_{\rho F}$ and the necessary kernel condition.  A cochain map
$C\to C'$ induces a natural map between these kernels for every $L$.
Conversely any natural transformation is determined on the two finite
sums of disks in the copresentation, hence on its kernel.  This proves
full faithfulness on natural transformations as well as on objects.
\end{proof}

\paragraph{A2.3 (Proved, and the shorter manuscript criterion refuted).}
\label{A2.3}
Suppose $R$ is a principal ideal domain.  The kernel conditions of A2.2
are equivalent to the existence, for all $i\geq0$, of $\ker d^i$ and,
for every nonzero nonunit $a\in R$, of
\begin{align*}
 K_{i,a}(C)&=\ker\left[
 C^i\oplus C^{i+1}\longrightarrow C^{i+1}\oplus C^{i+2}\right],\\
 (x,y)&\longmapsto(d^ix-ay,d^{i+1}y).
\end{align*}
For a field only the first family is needed.  For any nonfield PID
the first family alone is insufficient in general.  Here ``principal
ideal domain'' includes a field but excludes nonzero rings with zero
divisors.  For the broader meaning of ``principal ring'', A2.2 gives
the answer without a domain assumption; the PID reduction cannot be
used merely from principality of ideals.

\begin{proof}
We give the needed decomposition argument, so that the scope of the
reduction is explicit.  Write $Z_i=\ker\partial_i$ and
$B_i=\operatorname{im}\partial_{i+1}$.  Submodules of finite free modules
over a PID are finite free.  The sequence
$0\to Z_i\to L_i\to B_{i-1}\to0$ splits.  Choose such splittings.
Smith normal form for each inclusion $B_i\subset Z_i$ gives bases
$z_{ij}$ of $Z_i$ and $b_{ij}$ of $B_i$ with
$b_{ij}=a_{ij}z_{ij}$ for nonzero $a_{ij}$.  Lift the basis of
$B_{i-1}$ into the chosen complement of $Z_i$ in $L_i$.
The resulting bases decompose the entire complex, compatibly in all
degrees, into stalks and two-term complexes
\[
 A_i(a)=(R\xrightarrow{a}R)
       \quad\text{in homological degrees }i+1,i.
\]
Unit $a$ gives a disk, a zero differential gives two stalks, and the
remaining $a$ are precisely the stated nonzero nonunits.  For a stalk,
$T_C(S_i)=\ker d^i$; for a disk it is $C^i$; for $A_i(a)$ it is
$K_{i,a}(C)$.  Finite direct sums of these kernels give all the kernels
in A2.2.  Necessity follows by evaluating on these same complexes.

For a complete finite counterexample to the first-family-only claim,
take $R=\mathbf Z$ and let $\mathcal B$ be the full additive category
of finite free $\mathbf Z/4$-modules.  It is Karoubian: an idempotent
summand is a finite projective module over the local ring
$\mathbf Z/4$, hence free.  Take $C^1=\mathbf Z/4$ and $C^i=0$ for
$i\ne1$, all differentials zero.  Every $d^i$ has a kernel in
$\mathcal B$, but $K_{0,2}(C)$ would have to represent the kernel of
$2:\mathbf Z/4\to\mathbf Z/4$.  Evaluating its universal property at
the rank-one free object requires a group with two elements.  For a
free object of rank $r$, the group of maps from that generator has
$4^r$ elements; no nonnegative integer $r$ works.  Thus this kernel
does not exist.  The program exhausts the four elements, obtaining
$\{0,2\}$.  The same argument for a general nonfield PID uses
$R/(p^2)$ and the nonsplit submodule $pR/(p^2)$, where $p$ is an
irreducible.  That submodule is annihilated by $p$ and is nonzero,
whereas a nonzero free $R/(p^2)$-module is not annihilated by $p$.

Finally, over $R=\mathbf Z/4$, the three-term complex
$R\xrightarrow{2}R\xrightarrow{2}R$ is a complex, since $4=0$.
The PID proof does not apply: its boundaries are not projective.  It
cannot be a direct sum of stalks and two-term complexes of rank-one
free modules, since each of its three terms has rank one and both
successive differentials are nonzero.  A decomposition of the asserted
kind would require its middle rank-one term in two different summands.
This verifies the stated boundary of the elementary reduction.
\end{proof}

\paragraph{A2.4 (Proved; exact and admissible variants).}
\label{A2.4}
Let $(\mathcal E,\Sigma)$ be an $R$-linear Quillen exact category.
Restriction from exact functors
$\mathcal P_+(R)\to\mathcal E$ is fully faithful, with essential image
the nonnegative \emph{admissible resolutions}: complexes $C$ equipped
with objects $Z_i$ and conflations
\[
 0\longrightarrow Z_i\xrightarrow{s_i}C^i
    \xrightarrow{p_i}Z_{i+1}\longrightarrow0,
 \qquad d^i=s_{i+1}p_i\quad(i\geq0).
 \tag{A2.4a}
\]
In an abelian target with all short exact sequences this says exactly
$H^i(C)=0$ for $i>0$; $H^0(C)$ is unrestricted.

For strongly relatively left-exact functors, meaning those of A2.2
which also take source inflations to target inflations, the following
is the full criterion.  The kernels $T_C(L)$ of A2.2 must exist, and
for every degreewise split chain monomorphism $a:A\to B$ with
finite-free bounded quotient, the induced matrix map
$T_C(a):T_C(A)\to T_C(B)$ must be an inflation.  If $\mathcal E$ is
weakly idempotent complete, this second condition is equivalent to the
shorter requirement that each canonical inclusion
\[
 k_L:T_C(L)\longrightarrow\bigoplus_i L_i\otimes_R C^i
 \tag{A2.4b}
\]
be an inflation.  In particular this shorter test applies to
Karoubian exact categories.  Over a PID it suffices to test (A2.4b) on
the stalks and the two-term complexes of A2.3.

\begin{proof}
Full faithfulness follows from A2.2.  Exactness of $F$ applied to
$0\to S_i\to D_i\to S_{i+1}\to0$ proves necessity of (A2.4a), with
$Z_i=F(S_i)$.

For sufficiency we first work in an abelian target.  Given a complex
$C$ with no positive cohomology, set
$\mathcal H_L^q=\bigoplus_i L_i\otimes_R C^{i+q}$ with the usual
internal-Hom differential; thus $\mathcal H_L$ is
$\operatorname{Hom}^{\bullet}(L^\vee,C)$ and
$Z^0\mathcal H_L=T_C(L)$.  Its positive cohomology vanishes.  Indeed,
filter $L^\vee$ by its finitely many degrees.  The corresponding
successive quotients in the internal Hom are finite sums of shifts
$C[i]$ with $i\geq0$; their degree-$q$ cohomology is a finite sum of
$H^{q+i}(C)$, which is zero for $q>0$.  The long exact sequences
associated with this finite filtration prove the assertion.

A degreewise split sequence $0\to A\to B\to E\to0$ gives a
degreewise split short exact sequence
$0\to\mathcal H_A\to\mathcal H_B\to\mathcal H_E\to0$.
The obstruction to lifting a degree-zero cycle from $\mathcal H_E$
lies in $H^1\mathcal H_A=0$.  More explicitly, the exact cycle
sequence is
\[
 0\to Z^0\mathcal H_A\to Z^0\mathcal H_B
       \to Z^0\mathcal H_E\to H^1\mathcal H_A.
\]
It follows that $T_C$ is exact.

For an exact target we use the exact, fully faithful,
extension-closed abelian embedding of B\"uhler, Theorem~A.1
\cite{BuehlerExact}.  This theorem is applied in a universe containing
the objects and morphisms in question.  Its construction is compatible
with the central $R$-actions.  We must still show that the abelian
kernels $T_C(L)$ lie in $\mathcal E$; the embedding alone does not
imply that.  We prove this by induction on the support of $L$.
For $L=S_i\otimes_RV$ the object is $Z_i\otimes_RV$; for a disk it
is $C^i\otimes_RV$.  In general let $p$ be the least degree with
$L_p\ne0$, let $V=L_p$, and let $L'$ be the quotient by the stalk
$S_p\otimes_RV$.  There is a pullback square of complexes
\[
 \begin{array}{ccc}
 L&\longrightarrow&D_p\otimes_RV\\
 \downarrow&&\downarrow\\
 L'&\longrightarrow&S_{p+1}\otimes_RV.
 \end{array}
\]
The lower map is induced by $\partial_{p+1}$; the upper map has
components $1$ in degree $p$ and $\partial_{p+1}$ in degree $p+1$.
This verifies the square directly in every degree.  On applying
the abelian $T_C$ construction, the right vertical map becomes the
deflation $C^p\otimes_RV\to Z_{p+1}\otimes_RV$ from (A2.4a).
Its pullback along $T_C(L')\to Z_{p+1}\otimes_RV$ exists in
$\mathcal E$ by the Quillen pullback axiom, and represents $T_C(L)$.
The induction terminates because the support is finite.  The abelian
exactness already proved, together with reflection of exactness by
the embedding, now proves exactness in $\mathcal E$.

For the strong left-exact version, the criterion involving all
$T_C(a)$ follows from A2.2 and precisely specifies the extra
admissibility requirement using matrices, without presupposing any
extension functor.  Necessity of (A2.4b) follows by taking $a=j_L$
from A2.2.  Conversely suppose these inclusions are inflations and
$\mathcal E$ is weakly idempotent complete.  For a source inflation
$a:A\to B$, the map $I(a)$ is a split monomorphism between sums of
disks, because every $a_i$ splits.  Hence $T_C(I(a))$ is a split
monomorphism, and
\[
 k_B T_C(a)=T_C(I(a))k_A
\]
is an inflation.  Cancellation of inflations in a weakly idempotent
complete exact category, the dual of B\"uhler's Proposition~7.6,
implies that $T_C(a)$ is an inflation.  This proves the shorter test.
For a PID, the decomposition in A2.3 makes $k_L$ a finite direct sum
of the tested kernel inclusions and disk inclusions, up to invertible
changes of basis, so the final assertion follows.
\end{proof}

\paragraph{A2.5 (Refuted; the printed weak exactness axioms are insufficient).}
\label{A2.5}
The two axioms Sex~1--2 of 137:49, without the Quillen closure axioms,
do not imply the assertion that admissible resolutions are the
essential image of exact functors.  This remains false when the
underlying additive category is abelian.  The replacement under the
full exact-category hypotheses is A2.4.

\begin{proof}
Let $R=\mathbf Z$ and let $\mathcal E$ be the abelian category of
finite abelian groups.  Let $\Sigma$ consist of all sequences
isomorphic to finite direct sums of split short exact sequences and
copies of
\[
 e:\quad 0\longrightarrow\mathbf Z/4\xrightarrow{\,4\,}
 \mathbf Z/16\longrightarrow\mathbf Z/4\longrightarrow0.
\]
Every such sequence is a kernel--cokernel pair, and the class is
closed under isomorphisms and finite direct sums and contains all
split sequences.  Thus Sex~1--2 hold.  The complex
$C^0=\mathbf Z/16\to C^1=\mathbf Z/4$, with differential reduction
modulo four and all other terms zero, is a $\Sigma$-resolution:
its only nontrivial cycle sequence is $e$.

Suppose an exact functor extended $C$.  By A2.2 it would have to be
$T_C$.  For the degreewise split sequence
$0\to S_0\to A_0(2)\to S_1\to0$, its middle value is
\[
 G=\{(x,y)\in\mathbf Z/16\oplus\mathbf Z/4:
                      x\bmod4=2y\bmod4\}.
\]
This group has order sixteen and exponent eight: $(2,1)$ has order
eight, and every first coordinate is even.  The kernel of the
projection $G\to\mathbf Z/4$ is $\{(4t,0)\}$, a copy of
$\mathbf Z/4$, and the projection is surjective.  Therefore the
resulting sequence has both ends $\mathbf Z/4$ and middle group of
exponent eight.  It is not in $\Sigma$.  Indeed, a sequence of
$\Sigma$ with those ends can contain at most one copy of $e$;
if it contains one, all other summands have zero ends and its
middle exponent is sixteen.  If it contains none it is split,
with middle exponent four.  Both possibilities contradict exponent
eight.  The certificate exhausts all sixty-four pairs in the
definition of $G$, finds sixteen admissible pairs, and computes
their orders.  Thus the counterexample is finite and checks every
axiom being used.
\end{proof}

\paragraph{A2.6 (Proved; the valid Frobenius replacement for L3).}
\label{A2.6}
For every commutative $R$, the category
$\mathcal P_{\mathbf Z}(R)=\operatorname{Ch}^b
(\operatorname{Free}_{\mathrm{fg}}R)$ of complexes in all integer
degrees, with degreewise split sequences, is Frobenius.  Its
projective-injective objects are exactly its contractible complexes,
and its stable category is
$K^b(\operatorname{Free}_{\mathrm{fg}}R)$, the category of bounded
finite-free complexes and chain maps modulo chain homotopy.
The statement concerns finite-free complexes themselves; over a
general ring it does not assert that every cycle module of a
contractible complex is free.

\begin{proof}
The disk argument in A2.1 applies in every integer degree.  The
embedding $j_L:L\to I(L)$ from A2.2 provides enough injectives;
shifting the corresponding cone construction provides enough
projectives, with degreewise split maps and bounded finite-free
kernels and quotients.  Each $I(L)$ is contractible.  If $L$ is
injective, the embedding splits as a map of complexes, so $L$ is a
retract of a contractible complex and is contractible; the dual
argument applies to projectives.  Conversely, if $L$ is
contractible, any contraction splits its cone embedding and cone
epimorphism, so it is a retract of disk sums and hence both
injective and projective.  Retracts preserve the lifting properties
used here, whether or not individual cycle modules are free.

A map factoring through a contractible complex is null-homotopic.
Conversely a null-homotopic map factors through the standard
contractible cone of the identity on its source: the chosen
homotopy supplies the second component of the factorization, and
the equation $f=\partial h+h\partial$ is exactly its chain-map
condition.  Thus the ideal used in forming the stable category is
exactly the null-homotopy ideal.  This gives the asserted category
isomorphism.  The cone triangles match the stable triangles,
as in B\"uhler, Example~13.8 and Theorem~13.7
\cite{BuehlerExact}; the latter explicitly identifies Happel's
theorem as Chapter~I, Theorem~2.6, p.~16.  This is a statement about
the two-sided enlargement, not an equivalence of that enlargement
with $\mathcal P_+(R)$.
\end{proof}

